\chapter{Admissible Operating States and Constraint Projection}

\section{Entering Part III}

The first two parts of this book have used the language of admissibility, repair, and coupled constraints repeatedly, but always informally, as a way of talking about what the reactor's subsystems are doing to one another. Part III makes that language precise. This is not an exercise in mathematical decoration. Without a precise vocabulary for describing coupled admissibility, claims like "confinement and heating cannot be optimized independently" remain suggestive rather than actionable, and a suggestive claim cannot be used to design a control system, evaluate a safety case, or diagnose why a plant has drifted toward failure. This chapter builds the two central concepts the rest of Part III depends on: the admissible operating state, and constraint projection, the process by which a system is actively kept within it.

\section{The State of the Reactor}

At any moment, the reactor as a whole occupies some point in a very high-dimensional space. Some of its coordinates describe the plasma: density, temperature, and current profile at every point across the confinement volume. Others describe the magnetic field configuration, the temperature of every structural component, the tritium inventory and its distribution across breeding, extraction, and storage, the accumulated radiation damage in every material, the status of every cooling loop, and, in a fuller accounting, the staffing level and experience distribution of the operating crew and the state of the supply chain feeding replacement components. Call this entire collection the reactor's state, understanding from the outset that it spans physical, chemical, and organizational quantities together, not physical quantities alone.

Not every point in this space is one the reactor can occupy without damage. Some regions correspond to configurations the reactor can sustain indefinitely, or at least for as long as its design intends. Others correspond to configurations that will, if reached, trigger a disruption, breach a material's safe operating temperature, deplete tritium faster than it can be replenished, or otherwise push some subsystem past a threshold from which recovery is difficult or impossible. The admissible region is the subset of the state space the reactor can occupy while remaining consistent with its own continued operation. Everything the rest of this chapter says depends on taking this region seriously as a genuinely high-dimensional, jointly defined object, rather than as a convenient shorthand for a list of independent, subsystem-by-subsystem limits.

\section{Why the Admissible Region Cannot Be Factored}

It is tempting to imagine the admissible region as simply the intersection of separate limits: an allowed range for plasma density, an allowed range for magnet temperature, an allowed range for tritium inventory, each specified independently and then combined. This picture is wrong in a specific and consequential way. The limits themselves are frequently conditional on each other. The maximum plasma density the confinement system can safely sustain depends on the current heating profile. The rate at which tritium can safely be drawn from inventory depends on the current extraction system's throughput, which depends on the blanket's temperature, which depends on the cooling system's performance. A configuration that would be safe if only plasma density were considered, and separately safe if only magnet temperature were considered, can still be jointly unsafe, because the specific combination pushes some coupled quantity, not visible from either variable alone, outside its own tolerance.

This is why the admissible region has to be treated as a single, jointly defined object rather than a product of independently specified boxes. In practice, this means safety analysis and control system design cannot proceed by checking each subsystem's limits separately and declaring the reactor safe if all individual checks pass. It requires reasoning about the region's actual shape, which in a system with as many coupled subsystems as a fusion reactor is not a simple box at all, but a complicated, curved surface whose boundary in one direction depends on where the state currently sits along several other directions.

\section{Constraint Projection}

Given that the reactor's state is constantly drifting, driven by plasma instabilities on microsecond timescales, thermal cycling on timescales of hours, and material degradation on timescales of years, some active process has to continuously pull the state back toward the interior of the admissible region whenever it drifts too close to the boundary. Call this process constraint projection: the operation, performed by control systems, operators, and slower institutional processes together, of taking the reactor's current trajectory and correcting it so that it remains within, or is returned to, the admissible region.

The term is chosen deliberately. In the mathematics of constrained systems, projection refers to mapping a point that has strayed outside some allowed set back onto the nearest point within it. This is a reasonably faithful description of what a plasma control system does when it detects an instability beginning to grow and adjusts the magnetic field or heating profile to suppress it: the plasma's trajectory, left uncorrected, would carry it outside the admissible region, and the control system's intervention is best understood as continuously projecting that trajectory back toward safety, not as a discrete event that happens only when something goes wrong.

Constraint projection operates at every timescale this book has discussed. On the shortest timescales, feedback control systems adjust magnetic fields many times per second to suppress growing plasma instabilities. On intermediate timescales, operators adjust heating profiles, cooling flow rates, and tritium extraction rates in response to trends that unfold over minutes to hours. On the longest timescales, planned component replacement, blanket refurbishment, and operator retraining are themselves acts of constraint projection, correcting the slow drift of the admissible region's own boundary as materials degrade and institutional knowledge would otherwise erode. What looks, from a distance, like three unrelated activities, real-time plasma control, operational adjustment, and long-term maintenance planning, is, from the perspective developed in this chapter, the same underlying process operating at different timescales: continuously projecting the reactor's state back into a region from which it would otherwise drift away.

\section{Projection Has a Cost}

Constraint projection is not free. Every correction a control system makes consumes power, wears mechanical actuators, and, at the level of plasma control specifically, can itself introduce disturbances if applied too aggressively or with insufficient precision. A control system tuned to project the plasma's trajectory back toward the center of the admissible region as forcefully as possible, minimizing any excursion toward the boundary, may consume so much power or introduce so much actuator wear that it undermines the plant's overall energy balance or shortens the operating life of the components doing the correcting. This is a genuine trade-off, not a problem to be solved away by better engineering: staying further from the boundary is generally safer but more costly to maintain, while operating closer to the boundary reduces that cost but narrows the margin available before an excursion becomes a genuine violation.

This trade-off recurs at every timescale constraint projection operates on. A blanket refurbishment schedule that replaces components well before they approach their damage limits is safer but more expensive and involves more downtime than a schedule that waits until components are closer to their actual limits before replacing them. An operating crew maintained at a size well above the observed minimum needed for safe operation provides more margin against illness, turnover, and unexpected events, but costs more to sustain. None of these trade-offs has a universally correct answer. What viability theory offers is not a formula that resolves them automatically, but a consistent way of describing what is actually being traded: proximity to the boundary of the admissible region, in exchange for the cost of the projection needed to stay away from it.

\section{Toward Failure Cascades}

Constraint projection, described so far, has been treated as though it always succeeds, continuously pulling the reactor's state back into the admissible region no matter how far it drifts. Real systems do not have this guarantee. Projection has limits: a control system can only correct a plasma instability so fast, a cooling system can only remove heat so quickly, an operating crew can only respond to so many simultaneous demands. When the rate at which the reactor's state drifts toward the boundary of the admissible region exceeds the rate at which projection can correct it, the state crosses the boundary, and what happens next is the subject of the next chapter: how a single such crossing, in one subsystem, can cascade into failures across several others, on timescales ranging from microseconds to years, and what this implies for how a fusion reactor's overall safety and viability should actually be assessed.
